\section{Agents with Retrieval-Augmented Generation}
\label{sec:rag}

With the growing applications of LLMs, some circumstances require knowledge that is not contained in the training data set, for example, knowledge in highly professional domains or not publicly available. Even given the required datasets, the fine-tuning or re-training of the LLMs is still expensive. Accordingly, \textit{retrieval-augmented generation} (RAG), an innovative approach that aims to boost the power of LLMs in customized knowledge domain~\citep{rag_survey1,rag_survey2}, is gaining increasing attention in the literature.

The methodology of RAG can be considered as inserting a pre-processing step into the common utilization pipeline of LLMs. That is, given a collection of documents that contains needed knowledge, a similarity-based index is built, and the original user input is zipped with the most relevant pieces of information and converted into prompts, then fed to the LLMs. Therefore, the methodology of RAG involves multiple phases, that is, the collection of documents that contain the necessary information, the segmentation of the documents, the indexing of the segments (a.k.a. chunks or nodes), the similarity-based index retrieval, the fusion of the original query (i.e. user input) and retrieved results, the composing of prompts, and lastly, generation of reasonable responses from the LLM based on the informative prompts. 

In short, RAG embraces both the power of information retrieval and the generative capabilities of LLMs, and provides enhanced LLM service with customized domain knowledge at low cost. Meanwhile, assisted by RAG, the hallucination could be avoided and the factual accuracy could be significantly improved. 

As a developer-oriented multi-agent platform, \ours provides comprehensive RAG support for multi-agent applications. 
Given popular RAG frameworks such as \eg LlamaIndex~\citep{llamaindex}, LangChain~\citep{langchain}, \etc, \ours is designed with highly flexible abstracted processes to be compatible with those frameworks. In what follows, we introduce several key features of \ours RAG. 

% ===================== Save For Backup (5/10) =====================
% \zitao{
% Retrieval-augmented generation (RAG) is a methodology that combines the aspects of both information retrieval algorithms and generative models to enhance the performance and capabilities of natural language models. The workflow of RAG involves several steps. An index is first established for the loaded external data (which are split into pieces, a.k.a. chunks or nodes), therefore, closely related information pieces may admit smaller distances in terms of certain measures in vector space. Then, the user input is taken as a query, and the retriever searches the index (i.e. knowledge base) for information pieces with high similarities to the query. These retrieved information pieces are provided as context to compose a prompt and fed to the language generative models. With the in-context learning capability, the application of RAG would benefit LLMs from the following aspects. 1) LLMs can provide more reliable answers without requiring the question-related knowledge to be contained in the training dataset. 2) LLMs can generate answers that reflect the latest circumstances by frequently updating the knowledge base. 3) LLMs' hallucinations can be significantly reduced for specific knowledge domains if corresponding knowledge bases are loaded. 4) Customized LLM services can be provided without compromising data privacy, as no other party has direct access to the private knowledge base.
% }

% \zitao{
% In practice, RAG can be an important module for LLM-based agents. As a technique, RAG is closely related to the agent and multi-agent applications in the following aspects. On the application level, many multi-agent applications rely on the local knowledge base, such as a group of assistants providing Q\&A service based on multiple private datasets. From the capability of a more general single agent, RAG can improve some functionality, such as selecting the most appropriate tools from a large pool given a task description and detailed descriptions of the available tools. Besides the tool usage, when an agent has a long memory, RAG is also one of the default strategies to help agents extract the most relevant pieces to ensure the generated answers fit in the context. 
% }

% \paragraph{End-to-End JSON configurable process.}
% Different data preprocessing, such as operators splitting the long document, can lead to huge performance differences in the RAG technique. Because the multi-agent framework may have a wide range of applications, and the agents may need to work with diverse knowledge bases (e.g., PDF files, Python code base, or short messages). This demand requires that the RAG module in AgentScope should be flexible enough to work with different data preprocessing. As there are already many popular and highly-recognized RAG packages, such as LlamaIndex, AgentScope is designed to abstract the RAG process to be compatible with those existing packages. 
% Furthermore, AgentScope RAG modules are designed with the philosophy that most developers can use these modules and customize their own functions without writing Python code but indicate which available operator (e.g., CodeSplitter in LlamIndex) to pre-process the documents. This is achieved by modulizing the RAG into relatively independent submodules, including data loading, data processing, indexing, and retrieval. While the beginner can initialize the RAG modules with default functionalities in a few lines of code or config, the experienced developer can migrate or compose their RAG module as flexibly as they use packages initially.

\paragraph{Configurations in One-Stop}
Due to the complexity of the working pipeline, the configuration of RAG services is highly convoluted and often headache-some for users. While the RAG service provided by \ours is comprehensive and also involves multi-agent workflow, \ours provides a simple one-stop configuration solution by using a single \texttt{.json} file to group all RAG-related configurations. 

With this highly systematized configuration interface, users only need to focus on constructing the workflow, without being distracted by the repetitious configurations. For example, the RAG-empowered agents may involve a wide collection of knowledge bases that need to be configured in detail. With this ``One-Stop'' feature, the corresponding adjustments of the modules (which may lead to different performances) are integrated as the editing over simply one single file. Moreover, this solution also naturally adapts to the \ours Workstation, in which the dialog-box-based configuration can be easily exported to executable files and later loaded in Python programs.

% \paragraph{Knowledge-oriented Data Managements}
% When RAG is applied to memory management, it is required to be able to dynamically change the stored data: adding new memories, removing outdated or conflicting memories, and refreshing the memory indexing. Motivated by such demand, the RAG modules in AgentScope inherit the capability of the modules to update the data. 
% Moreover, we recognize the potential for applications where multiple agents share the same knowledge base, such as in a customer service system. In such scenarios, it is impractical to have each agent initialize their RAG indexing independently. Instead, it is more efficient to allow them to share the embedding and indexing of the knowledge base, thereby avoiding repeated calls to embedding models and reassuring the audience about the practicality of this approach.
% However, retrieval can be decoupled from the indexing phase. For example, although two agents share the same knowledge base with indexing and embeddings, they may have different strategies in query-rewrite or post-retrieval ranking. Therefore, we allow each agent to have their own retrievers (i.e., retrieval strategies) while building those from the same set of indexing and embeddings.

\paragraph{Knowledge-oriented Data Managements}
The application of RAG in multi-agent circumstances is more complicated compared with the application on a single agent. For example, for a single agent, one can directly encapsulate the needed knowledge to the agent. Therefore, the initialization of each RAG agent involves the whole pipeline of conversion from the original documents to vector-stored indexes with retrievers. However, in multi-agent applications, it is natural for agents to share knowledge, such that repeatedly executing index computation for each agent is needless. Therefore, \ours introduces the notion of \texttt{knowledge banks}.

Knowledge banks can be considered as a collection of knowledge containers, where the smallest manageable unit is a customized object (which will be referred to as a ``RAG object'' in the following context). The workflow starts with initializing the knowledge bank, which mainly relies on the information contained in the \texttt{.json} configuration file. The information includes the directory and extensions (such as \texttt{.py} or \texttt{.md}) of documents, the granularity and choice of segmentation tools (e.g. the splitters in \texttt{Llama-Index}) for the documents, and choice of model for indexing. After the initialization, the computed results are persisted to the designated directory for later use and we also obtain a knowledge bank consisting of RAG objects, each marked with a unique \texttt{knowledge\_id}, associated with the index of the corresponding documents, an information retriever, and other attributes. Note that \ours permits each RAG agent to load with more than one RAG object.

% \paragraph{Agents working with RAG}
% Compared with using the RAG techniques directly, using RAG in \ours provides additional probabilities for improving the user experience of RAG.
% \begin{itemize}
%     \item 
%     The parallelizable execution of the agents can accelerate the preprocessing process.
%     \item 
%     Agents, when tasked with generating answers using RAG, can be programmed to rewrite queries based on their internal memory settings. This approach makes the retrieval process more customizable and agent-oriented. 
%     \item 
%     With the multi-agent framework, RAG's process can be more flexible. When generating answers based on a knowledge base, the knowledge can be either aggregated as a complete set for retrieval or divided and assigned to different agents for retrieval but summarized the answers of those agents at last.
% \end{itemize}

\paragraph{Agents with RAG}

The application of agents with RAG in \ours is very simple. For example, we first need to initialize a \texttt{KnowledgeBank} with some RAG framework, e.g. LlamaIndex, and all the documents. Then, we configure an RAG agent and load it with the knowledge bank. After that, the initialization is completed and we can use the RAG agent like any other agent in \ours. It is worth noting that if \texttt{KnowldgeBank} is obtained with LlamaIndex framework, then we need to use \texttt{LlamaIndexAgent} (inherited from \texttt{RAGAgentBase}). The readers may refer to Section~\ref{sec:copilot} for a concrete application sample, which implements a copilot for \ours using our RAG agents. Overall, the key features of RAG agents are summarized as follows:
\begin{itemize}
    \item The RAG agent is permitted to load several RAG objects (i.e. any subset of the knowledge bank). One can choose to load the original RAG objects from the knowledge bank (in such case, the modification to an object may affect all the agents who use it) or a copy of it.
    \item While agents are initialized with a \texttt{KnowledgeBank} object, it is permitted for the agents to update knowledge in time. The operations include inserting, deleting, or replacing knowledge pieces. Moreover, we provide a solution by monitoring certain directories and keeping the RAG object updated with the contents in the directories.
    \item The fusion mechanism of the retrieved results from multiple RAG objects is fully customizable. For example, since knowledge may be of different importance or trustworthiness, the agent can set weights for information retrieved from different RAG objects for subsequent processes. 
    \item RAGs agents are permitted to recompose the query in configurable repeats and conduct multiple queries for more comprehensive answers.
\end{itemize}